\color{unidarkgreen}{\subsection{Логика деления стороны ВН трансформатора}}
\color{black}

\begin{enumerate}

\item
Ввод в работу логики деления стороны ВН трансформатора от ТЗНП осуществляется с помощью уставки \textbf{«ЛД ВН от ТЗНП»}. 

\item
Алгоритм работы логики деления стороны ВН трансформатора приведен на рисунке \ref{pict:LD_VN}, уставки функции приведены в таблице \ref{app_set:ldvn}.

\vspace{1mm}
\begin{figure}[H]
\centering
\includegraphics[width=0.8\textwidth,height=0.8\textheight,keepaspectratio]{\fbpath/07.01 РЗТ/021.1791 ЛД ВН/img/LDVN.pdf}
\captionof{figure}{Функциональная схема ЛД стороны ВН }\label{pict:LD_VN}
\end{figure}

\item
Оперативный ввод/вывод ЛД высокой стороны трансформатора производится с помощью БУ (блока управления функциями) \textit{«Вывод ЛД ВН»}. 

\item
Первым действием от пуска ТЗНП ВН через выдержку времени \textbf{«Тср ТЗНП откл СВ(ШСВ) ВН»} логика деления формирует сигналы предварительного деления с действием на отключение СВ (ШСВ) ВН.

\item
Вторым действием от пуска ТЗНП ВН через выдержку времени \textbf{«Тср ТЗНП откл В ВН»} логика деления формирует сигналы на отключение выключателей высокой стороны трансформатора.

\end{enumerate} 